% ===== PRÉAMBULE CANONIQUE — à coller VERBATIM en tête de CHAQUE .tex =====
\documentclass[11pt,a4paper]{article}
\usepackage[utf8]{inputenc}\usepackage[T1]{fontenc}\usepackage{lmodern}
\usepackage[french]{babel}
\usepackage{amsmath,amssymb,mathtools}\usepackage{icomma}
\usepackage[version=4]{mhchem}          % \ce{...} : équations & formules
\usepackage{siunitx}\sisetup{locale=FR} % \qty{}{}, \si{}, \ang{}
\usepackage{chemfig}                    % \chemfig{...} : structures développées
\usepackage{xcolor}\usepackage{colortbl}
\usepackage{tikz}\usepackage{pgfplots}\pgfplotsset{compat=1.18}
\usepackage[most]{tcolorbox}\usepackage{enumitem}\usepackage{array,booktabs}
\usepackage[a4paper,margin=2.2cm]{geometry}
\usetikzlibrary{arrows.meta,calc,angles,quotes,positioning,patterns,backgrounds,shapes.geometric}
\definecolor{primary}{RGB}{222,49,99}\definecolor{primarydark}{RGB}{150,22,55}
\definecolor{defcol}{RGB}{0,114,178}\definecolor{methcol}{RGB}{52,92,148}
\definecolor{excol}{RGB}{0,135,90}\definecolor{attncol}{RGB}{200,40,55}
\tcbset{boxbase/.style={enhanced,breakable,boxrule=0.9pt,arc=2.5pt,
  left=3mm,right=3mm,top=2.3mm,bottom=2.3mm,fonttitle=\bfseries,coltitle=white}}
% --- boîtes TUTORIEL (noms EXACTS, ne pas renommer) ---
\newtcolorbox{cle}[1][]{boxbase,colback=defcol!6,colframe=defcol,
  title={\textsc{À retenir}\ifx\relax#1\relax\relax\else\textmd{ : \itshape #1}\fi}}
\newtcolorbox{methode}[1][]{boxbase,colback=methcol!7,colframe=methcol,sharp corners,
  title={\textsc{Méthode}\ifx\relax#1\relax\relax\else\textmd{ --- \itshape #1}\fi}}
\newtcolorbox{exemplebox}[1][]{boxbase,colback=excol!5,colframe=excol,
  title={\textsc{Exemple}\ifx\relax#1\relax\relax\else\textmd{ : \itshape #1}\fi}}
\newtcolorbox{piege}[1][]{boxbase,colback=attncol!6,colframe=attncol,
  title={\textsc{Piège}\ifx\relax#1\relax\relax\else\textmd{ : \itshape #1}\fi}}
\newtcolorbox{remarque}[1][]{boxbase,colback=black!4,colframe=black!45,
  title={\textsc{Remarque}\ifx\relax#1\relax\relax\else\textmd{ : \itshape #1}\fi}}
% --- boîtes EXERCICE / CORRIGÉ (noms EXACTS, auto-comptées) ---
\newtcolorbox[auto counter]{exo}[1][]{boxbase,colback=primary!4,colframe=primary,
  title={\textsc{Exercice}~\thetcbcounter\ifx\relax#1\relax\relax\else\textmd{ --- \itshape #1}\fi}}
\newtcolorbox[auto counter]{corrige}[1][]{boxbase,colback=excol!4,colframe=excol,
  title={\textsc{Corrigé}~\thetcbcounter\ifx\relax#1\relax\relax\else\textmd{ --- \itshape #1}\fi}}
\newcommand{\R}{\mathbb{R}}\newcommand{\N}{\mathbb{N}}\newcommand{\Z}{\mathbb{Z}}\newcommand{\ds}{\displaystyle}
% =========================================================================
\begin{document}
\begin{exo}[escalier complet : de la glace à la vapeur]
On veut la chaleur totale $Q$ pour transformer $\qty{50}{\gram}$ de glace prise à $\qty{-40}{\celsius}$ en vapeur portée à $\qty{120}{\celsius}$. Données : $c_{\text{glace}}=\num{2.060}$, $c_{\text{eau}}=\num{4.185}$, $c_{\text{vapeur}}=\num{1.85}$ (tous en $\si{\kilo\joule\per\kelvin\per\kilogram}$), $L_{\text{fus}}=\qty{334}{\kilo\joule\per\kilogram}$, $L_{\text{vap}}=\qty{2257}{\kilo\joule\per\kilogram}$.
\begin{center}
\begin{tikzpicture}[scale=1.0,every node/.style={font=\scriptsize}]
  \draw[->,line width=1.1pt,defcol] (0,0)--(2.1,0);
  \node[above] at (1.05,0) {$m\,c_{\text{glace}}\Delta T$};
  \node[below left] at (0,0) {$\qty{-40}{\celsius}$}; \node[below] at (0,-0.38) {glace};
  \draw[->,line width=1.1pt,attncol] (2.1,0)--(2.1,1.0);
  \node[right] at (2.1,0.5) {$m\,L_{\text{fus}}$}; \node[below] at (2.1,0) {$\qty{0}{\celsius}$};
  \draw[->,line width=1.1pt,defcol] (2.1,1.0)--(5.1,1.0);
  \node[above] at (3.6,1.0) {$m\,c_{\text{eau}}\Delta T$};
  \draw[->,line width=1.1pt,attncol] (5.1,1.0)--(5.1,2.0);
  \node[right] at (5.1,1.5) {$m\,L_{\text{vap}}$}; \node[below] at (5.1,1.0) {$\qty{100}{\celsius}$};
  \draw[->,line width=1.1pt,defcol] (5.1,2.0)--(7.5,2.0);
  \node[above] at (6.3,2.0) {$m\,c_{\text{vap}}\Delta T$};
  \node[above right] at (7.5,2.0) {$\qty{120}{\celsius}$}; \node[below] at (7.5,2.0) {vapeur};
\end{tikzpicture}
\end{center}
\begin{enumerate}[label=\alph*)]
  \item Calcule les cinq contributions séparément.
  \item En déduis $Q$ total.
  \item Quelle étape domine le bilan ? Justifie.
\end{enumerate}
\end{exo}

\begin{exo}[dissolution : compter aussi le calorimètre]
Un calorimètre en verre de masse $\qty{0.30}{\kilogram}$ contient $\qty{0.25}{\kilogram}$ d'eau à $\qty{21.0}{\celsius}$. On y dissout $\qty{0.10}{\mol}$ d'un sel ; la température monte à $\qty{24.0}{\celsius}$. On suppose que \emph{seule la moitié} du verre s'échauffe. Données : $c_{\text{verre}}=\qty{0.8368}{\kilo\joule\per\kelvin\per\kilogram}$, $c_{\text{eau}}=\qty{4.185}{\kilo\joule\per\kelvin\per\kilogram}$.
\begin{enumerate}[label=\alph*)]
  \item Calcule la chaleur captée par l'eau.
  \item Calcule la chaleur captée par le verre.
  \item En déduis l'enthalpie molaire de dissolution $\Delta H$.
\end{enumerate}
\end{exo}

\begin{exo}[une barre de fer froide : congélation partielle]
Un calorimètre adiabatique (capacité négligeable) contient $\qty{0.25}{\kilogram}$ d'eau à $\qty{4}{\celsius}$. On y plonge une barre de fer de $\qty{0.60}{\kilogram}$ sortant d'un congélateur à $\qty{-35}{\celsius}$. Données : $c_{\text{eau}}=\qty{4}{\kilo\joule\per\kelvin\per\kilogram}$, $c_{\text{fer}}=\qty{0.5}{\kilo\joule\per\kelvin\per\kilogram}$, $L_{\text{fus}}=\qty{334}{\kilo\joule\per\kilogram}$.
\begin{enumerate}[label=\alph*)]
  \item Quelle énergie faut-il fournir au fer pour l'amener de $\qty{-35}{\celsius}$ à $\qty{0}{\celsius}$ ?
  \item Quelle énergie l'eau peut-elle céder en se refroidissant de $\qty{4}{\celsius}$ à $\qty{0}{\celsius}$ ?
  \item Montre qu'une partie seulement de l'eau gèle.
  \item Calcule la masse d'eau gelée.
  \item Quelle est la température finale d'équilibre ?
\end{enumerate}
\end{exo}

\begin{exo}[vrai ou faux — calorimétrie]
Indique si chaque affirmation est \textbf{vraie} ou \textbf{fausse}, et justifie brièvement.
\begin{enumerate}[label=\alph*)]
  \item « Pendant un changement d'état, la température du corps reste constante. »
  \item « Une variation de $\qty{40}{\kelvin}$ correspond à une variation de $\qty{40}{\celsius}$. »
  \item « Dans un calorimètre adiabatique, la somme des chaleurs échangées entre les corps présents est nulle. »
  \item « $\qty{4.185}{\kilo\joule}$ valent $\qty{41.85}{\joule}$. »
\end{enumerate}
\end{exo}

\end{document}
